\section{Preliminary experiments}\label{app: prel exp}

\subsection{Noise structure}
We generate the labels with a centered parabola (see Section~\ref{app/ sec: datagen}), so that roughly half of them are $1$ and half $-1$ in the interval $[0,10]$. Then, we inspect what happens if we use a Bernoullian noise instead of a Gaussian one, i.e. when data is generated as 
\[
y \leftarrow sign(h^*(x))b
\]
for $b\sim\mathrm{Bernoulli}(p)$. Notice that this noise flips the labels with a probability of $1-p$. 
While from the first experiments it seemed that the algorithm achieved a higher test error, we postulate that this may be due to the portion of misclassified inputs, rather than to the noise structure itself. Hence, fixing $p = 0.8$, we look for the value of $\sigma$ that achieves roughly the same portion of mislcassified inputs ($20\%$). Figure~\ref{app/fig: sigmastar parabola} shows that this happens for $\sigma = 5.5$, and Figures~\ref{app/fig: parabola bernoulli} and~\ref{app/fig: parabola gauss} show that the two noise structures then yield essentially the same train and test errors. 

\begin{figure}[h!]
  \centering
  \includegraphics[width=0.8\textwidth]{../../outputs/recap/app_sigmastar_parabola.png}
  \caption{Average fraction of misclassified inputs against $\sigma$, over $100$ samples of size $n = 50$, $h^*$ = centered parabola}
  \label{app/fig: sigmastar parabola}
\end{figure}

\begin{figure}[h!]
  \centering
  \begin{minipage}{0.48\textwidth}
    \centering
    \includegraphics[width=\textwidth]{../../outputs/recap/app_parabola_bernoulli.png}
    \caption{$n = 50, N = 10,000, n_{\text{test}} = 200,000$, Bernoulli noise with $p = 0.8$}
    \label{app/fig: parabola bernoulli}
  \end{minipage}\hfill
  \begin{minipage}{0.48\textwidth}
    \centering
    \includegraphics[width=\textwidth]{../../outputs/recap/app_parabola_gauss.png}
    \caption{$n = 50, N = 10,000, n_{\text{test}} = 200,000$, Gaussian noise with $\sigma = 5.5$}
    \label{app/fig: parabola gauss}
  \end{minipage}
\end{figure}

\subsection{Data-generating function}\label{app/ sec: datagen}
Our first choice of $h^*$ is a parabola centered horizontally at $5$, the midpoint of the support $\mathcal X$, and vertically so that roughly half of the points fall above the $x$-axis and half below. I.e.,
\[
h^*(x) = (x-5)^2 -6.25
\]

\begin{figure}[h!]
  \centering
  \includegraphics[width=0.7\textwidth]{../../outputs/recap/app_explore_parabola.png}
  \caption{A sample of size $n = 50$ generated with the centered parabola, for $\sigma \in \{10^{-4}, 1, 5.5, 10\}$}
  \label{app/fig: explore data parabola}
\end{figure}

Figure~\ref{app/fig: explore data parabola} shows the data generated with the centered parabola, for different amounts of noise.

Because in theory an algorithm could learn such a classification pattern as soon as it sees $3$ points (one point for each region), we investigate what happens when the underlying function is more complex. We study a sinusoid with varying frequency,
\[
    h^*(x) = 5\, \sin\left(\frac{\pi \,( m  +1)}{10}\, x\right)
\]
We multiply the sine by $5$ to allow for enough height to properly inspect the data, while the argument of the sine is such that $m$ is the number of times $h^*$ crosses the $x$-axis in the open interval $(0,10)$, yielding $m+1$ contiguous blocks of $1s$ and $-1s$. 

\begin{figure}[h!]
  \centering
  \begin{minipage}{0.32\textwidth}
    \centering
    \includegraphics[width=\textwidth]{../../outputs/recap/app_sine_m10.png}
    \caption{$m = 10$}
    \label{app/fig: sine m10}
  \end{minipage}\hfill
  \begin{minipage}{0.32\textwidth}
    \centering
    \includegraphics[width=\textwidth]{../../outputs/recap/app_sine_m100.png}
    \caption{$m = 100$}
    \label{app/fig: sine m100}
  \end{minipage}\hfill
  \begin{minipage}{0.32\textwidth}
    \centering
    \includegraphics[width=\textwidth]{../../outputs/recap/app_sine_m1000.png}
    \caption{$m = 1000$}
    \label{app/fig: sine m1000}
  \end{minipage}
\end{figure}

\begin{figure}[h!]
  \centering
  \includegraphics[width=0.8\textwidth]{../../outputs/recap/app_ratio_sweep.png}
  \caption{$n = 50, N = 10,000, n_{\text{test}} = 200,000$, $\sigma = 3$, train and test error for different ratios $r = m/n$}
  \label{app/fig: ratio sweep}
\end{figure}

We explore how a random sample and the performance of the algorithm change for different sine frequencies. In particular, we expect that for a high frequency the algorithm will not learn, as our datasize is $n=50$.
In Figures~\ref{app/fig: sine m10}--\ref{app/fig: sine m1000}, we test a function that osciallets $10, 100$, and $1000$ times with a noise level of $\sigma = 3$ (we reduce it compared to $\sigma = 5$, having increased the complexity of the function). We notice that for $m=1000$ the algorithm fails to learn, and the test error stays stable at $0.5$. This is because the labels change so frequently with respect to the number of points in the training dataset $n$ that it is as if we had random labels. Hence, we investigate the relation between the complexity of the underlying function $m$ and the data size $n$. We find that for $\frac{m}{n}>5$, roughly, the algorithm ceases to learn (Figure~\ref{app/fig: ratio sweep}). In the following experiments, we fix $m=10$ and $n=50$, which is very far from the threshold.

\begin{figure}[h!]
  \centering
  \includegraphics[width=0.8\textwidth]{../../outputs/recap/app_misclassified_sine.png}
  \caption{Fraction of misclassified inputs over $100$ samples of size $n = 50$, $h^*$ = sine with $m = 10$; dashed lines are the averages}
  \label{app/fig: misclassified sine}
\end{figure}

\begin{figure}[h!]
  \centering
  \includegraphics[width=0.8\textwidth]{../../outputs/recap/app_sigmastar_sine.png}
  \caption{Average fraction of misclassified inputs against $\sigma$, over $100$ samples of size $n = 50$, $h^*$ = sine with $m = 10$}
  \label{app/fig: sigmastar sine}
\end{figure}

We then verify whether a Gaussian noise with $\sigma = 5.5$ achieves the same portion of misclassified inputs on average as a Bernoullian noise with $p=0.8$, as was the case for data generated with a parabola. Figure~\ref{app/fig: misclassified sine} shows that this is not the case on average ($29\%$ against $20\%$), while Figure~\ref{app/fig: sigmastar sine} shows that the desired value is $\sigma = 3.5$.

\begin{figure}[h!]
  \centering
  \includegraphics[width=0.8\textwidth]{../../outputs/recap/app_simple_function_sigma1.png}
  \caption{$N = 10,000, n_{\text{test}} = 200,000, \sigma = 1$, $h^*$ = sine with $m = 2$: test error at depth $2$ and $3$ against $n$, averaged over $100$ samples}
  \label{app/fig: simple function sigma1}
\end{figure}

\begin{figure}[h!]
  \centering
  \includegraphics[width=0.8\textwidth]{../../outputs/recap/app_simple_function.png}
  \caption{$N = 10,000, n_{\text{test}} = 200,000, \sigma = 3.5$, $h^*$ = sine with $m = 2$: test error at depth $2$ and $3$ against $n$, averaged over $100$ samples}
  \label{app/fig: simple function}
\end{figure}

Finally, we investigate the afore-mentioned intuition about how easy it is for the algorithm to learn a simple function. We take $m=2$ (which corresponds to a parabola, in terms of classification pattern), sample uniformly from $\X$, and look for the training set size $n$ for which the test error gets within $0.01$ of the Bayes risk on average, at depth $2$ and $3$.

We start with little noise, $\sigma = 1$, for which the Bayes risk is $\approx 0.05$ (Figure~\ref{app/fig: simple function sigma1}). Here the algorithm learns the function quickly: with $n = 50$ points the test error is already $\approx 0.07$, within $0.02$ of the Bayes risk, and it gets within $0.01$ with $n = 200$ at depth $2$ and $n = 500$ at depth $3$.

With our noise level, $\sigma = 3.5$, the Bayes risk is $\approx 0.20$ and $n = 1000$ training points are needed on average for the test error to get within $0.01$ of it, both at depth $2$ and $3$ (Figure~\ref{app/fig: simple function}). Hence, although two splits are enough to represent the function, with this amount of noise the algorithm needs many points to place them correctly.


\newpage
\section{Proofs}\label{app: proofs}
\begin{proof}[Proof of Theorem~\ref{thm: unionbound}]
  Apply Hoeffding Lemma and a union bound over the hypotheses set $\mathcal H$. Tail integrate to obtain a result in expected value. 
\end{proof}

\begin{proof}[Proof of Theorem~\ref{thm: rad bound}]
  Symmetrization technique. 
\end{proof}

\begin{proof}[Proof of Lemma~\ref{lemma: refined massart}]
  Pass to the supremum over {\it loss vectors\/} instead of functions $\ell\circ h$, and bound it with the sum, which makes the $N(D)$ term appear. Then, apply the exponential Hoeffding Lemma and optimize for the coefficient in the exponent. 
\end{proof}

\begin{proof}[Proof of Lemma~\ref{lemma: counting}]
  Given a dataset, its realizations through $\ell_h$ are less than or equal to those through $h$ because applying $\ell$ cannot increase the number of outputs, and the number of images of the functions in $\mathcal H$ are less than or equal to the number of functions $\card$. On the other hand, $N_D(\cH)\le 2^n$ because we have at most $2^n$ $n$-long sequences of $0s$ and $1s$.
\end{proof}

\begin{proof}[Proof of Theorem~\ref{thm: rademmassart bound}]
  Apply Theorem~\ref{thm: rad bound}, Lemma~\ref{lemma: refined massart} and Lemma~\ref{lemma: counting}. 
\end{proof}

\begin{proof}[Proof of Theorem~\ref{thm: bound emprisk}]
        Apply the relative deviation inequality (see Loz's file),
        \[
        \Pr\left[\exists h\in\cH : R(h)-\hat R(h) > \varepsilon \sqrt{R(h)}\right] \le 2N_{2n} (\cH) e^{-2\varepsilon^2/16}
        \]
        with the suitable $\varepsilon$ and then isolate $\hat R(\hat h)$ on the RHS.
        % I have not looked precisely at the proof of the rel dev ineq
\end{proof}

\begin{proof}[Proof of Corollary~\ref{coroll: bound bestrisk}]
         Set $h_0 = \argmin_{h\in\cH} R(h)$, so that $R(h_0) = R^*_{\cH}$. Then,
        \begin{equation}\label{eq: hatR h le hat R h0}
        \hat R(\hat h) \le \hat R(h_0).
        \end{equation}
        Since $h_0$ is independent of the data, $\expv_D[\hat R(h_0)] = R(h_0)$ and hence Bernstein gives, with probability at least $1-\delta/2$,
        \[
        \hat R(h_0) \le R^*_{\cH} + \sqrt{\frac{2R*_\cH \log (2/\delta)}{n}} + \frac{\log(2/\delta)}{3n}.
        \]
        Since $N_{2n}(\cH)\ge 1$ and thus $\frac{\log (2/\delta)}{n}\le \frac{\log(4N_{2n}(\cH)/\delta)}{n}$,
        \begin{equation}\label{eq: bound hatR h0}
        \hat R(h_0) \le \left(\sqrt{R^*_\cH} + \sqrt{\frac{4\log(N_{2n}(\cH)/\delta)}{n}}\right)^2
        \end{equation}
        Substitute~(\ref{eq: hatR h le hat R h0}) into~(\ref{eq: bound hatR h0}), and apply the bound to $\hat R(\hat h)$ in the RHS of Theorem~\ref{thm: bound emprisk} at confidence $1-\delta/2$. This yields the desired bound, which holds with probability at least $1-\delta$ by a union bound on the events ((\ref{eq: bound hatR h0}) and Theorem~\ref{thm: bound emprisk}) that we require to hold simultaneously.
\end{proof}

\begin{proof}[Proof of Corollary~\ref{coroll: bound noiseandappr}]
  Substitute $R^*_\cH = \sigma^2 + A_\cH$ into Corollary~\ref{coroll: bound bestrisk}.
\end{proof}

\begin{proof}[Proof of Theorem~\ref{thm: bound noiseandappr}]
    Let $\text{gap} = (R-\hat R(\hat h))$ and $C = \sqrt{\frac{\log(4 N_n(\cH)/\delta)}{n}}$. Corollary~\ref{cor: best risk} gives
\begin{align*}
&\Pr\left[\text{gap} \le \sqrt{R^*_\cH}\,C + 20C^2\right] \ge 1-\delta \\
\implies &\Pr\left[\text{gap} > \sqrt{R^*_\cH}\,C + 20C^2\right] \le \delta
\end{align*}
Set $\varepsilon = \sqrt{R^*_\cH}\,C + 20C^2$. Let us now find $\delta(\varepsilon)$:
\[
20C^2 + \sqrt{R^*_\cH}\,C - \varepsilon = 0 \implies C = \frac{-\sqrt{R^*_\cH} + \sqrt{R^*_\cH + 80\varepsilon}}{40}
\]
So we have
\begin{align*}
\frac{\log(4N_{2n}(\cH)/\delta)}{n} &= \left(\frac{\sqrt{R^*_\cH + 80\varepsilon} - \sqrt{R^*_\cH}}{40}\right)^2 \\
\frac{4N_{2n}(\cH)}{\delta} &= \exp\left(n\left(\frac{\sqrt{R^*_\cH + 80\varepsilon} - \sqrt{R^*_\cH}}{40}\right)^2\right)
\end{align*}
and hence
\[
\delta(\varepsilon) = 4N_{2n}(\cH)\exp\left(-\frac{n}{1600}\left(\sqrt{R^*_\cH + 80\varepsilon} - \sqrt{R^*_\cH}\right)^2\right).
\]
Using that $\expv[X] = \int_0^\infty \Pr[X \ge \varepsilon]\,d\varepsilon$ for $X \ge 0$, applied to $\text{gap}^+ \coloneq \max\{\text{gap}, 0\}$,
\[
\expv_D[\text{gap}] \le \expv_D[\text{gap}^+] = \int_0^\infty \Pr[\text{gap} \ge \varepsilon]\,d\varepsilon \le \int_0^\infty \min\{1, \delta(\varepsilon)\}\,d\varepsilon.
\]
By a shift of perspective, looking at the integral as a sum of horizontal stripes instead of vertical stripes in the graph of $\varepsilon \mapsto \delta(\varepsilon)$, we obtain
\[
\int_0^\infty \min\{1, \delta(\varepsilon)\}\,d\varepsilon = \int_0^1 \delta^{-1}(t)\,dt = \int_0^1 \varepsilon(\delta)\,d\delta
\]
Hence, writing $\varepsilon(\delta) = \sqrt{R^*_\cH}\,C + 20C^2$ explicitly,
\[
\expv_D[\text{gap}] \le \int_0^1 \varepsilon(\delta)\,d\delta = \sqrt{\frac{R^*_\cH}{n}}\int_0^1 \sqrt{\log\frac{4N_{2n}(\cH)}{\delta}}\,d\delta + \frac{20}{n}\int_0^1 \log\frac{4N_{2n}(\cH)}{\delta}\,d\delta.
\]
The logarithm is the only part depending on $\delta$, and its integral is
\[
\int_0^1 \log\frac{4N_{2n}(\cH)}{\delta}\,d\delta = \log(4N_{2n}(\cH)) - \int_0^1 \log\delta\,d\delta = \log(4N_{2n}(\cH)) + 1.
\]
Applying Jensen's inequality to the square root term,
\[
\int_0^1 \sqrt{\log\frac{4N_{2n}(\cH)}{\delta}}\,d\delta \le \sqrt{\int_0^1 \log\frac{4N_{2n}(\cH)}{\delta}\,d\delta} = \sqrt{\log(4N_{2n}(\cH)) + 1},
\]
we obtain
\[
\expv_D[\text{gap}] \le \sqrt{\frac{R^*_\cH\left(\log(4N_{2n}(\cH)) + 1\right)}{n}} + 20\,\frac{\log(4N_{2n}(\cH)) + 1}{n}.
\]

Asymptotically,
\[
\expv_D\left[R(\hat h)-\hat R(\hat h)\right] \le O\left(
\sqrt{\frac{(A_{\mathcal H}+\sigma^2)\,\log N_{2n}(\mathcal H)}{n}} + \frac{\log N_{2n}(\mathcal H)}{n} \right)
\]
because, at the numerator,
\[
\log (4 N_{2n}(\mathcal H)) + 1 = \log 4 + \log N_{2n}(\mathcal H) + 1 = \log N_{2n}(\mathcal H) + O(1)
\]
and for large $x$, 
\[
\log x + O(1) \le C \log x
\] 
for some constant $C>0$. Plugging in the counting inequality from Lemma~\ref{lemma: counting} with $2n$, we obtain and plot the following bound
\[
\expv_D\left[R(\hat h)-\hat R(\hat h)\right] \le \min\left\{ O\left(
\sqrt{\frac{(A_{\mathcal H}+\sigma^2)\log \card}{n}} + \frac{\log\card}{n} \right), O\left(
\sqrt{A_{\mathcal H}+\sigma^2}+1  \right)\right\}
\]
Remark: recall that $N_{2n}(\cH)$ is the maximum number of loss vectors over the union of two datasets $D, D'$, or equivalently over a dataset with size $2n$. The counting lemma must be applied to $2n$-long loss vectors. Hence, we still have $N_{2n}(\cH) \le \vert \cH\vert$ because the number of functions is still the same, and the bound counting all the possible sequences becomes $N_{2n}(\cH) \le 2^{2n}$. 

\end{proof}

\begin{proof}[Proof of Lemma~\ref{lemma: growth}]
    $N_D(\cH) \le \Pi_{n}(\cH)$ because the number of distinct loss vector can at most decrease, given prediction vectors. $\Pi_{n}(\cH)\le \card$ because the realizations of functions in $h$ are at most as much as the functions in $h$ (equality when functions in $h$ are bijective). $\Pi_{n}(\cH)\le 2^n$ because we have at most $2^n$ sequences of $n$ $1$s and $-1$s.
\end{proof}

\newpage
\section{Additional results}\label{app: add res}

\begin{fact}[The empirical risk is a biased estimate of the true risk for the ERM]
lalala
  
\end{fact}

\begin{fact}[Bound on the Rademacher complexity of $\LL$]
lalalala

\end{fact}